% Reconstructed from the complete reviewed chat text after the workspace outage.
% This is a mathematical transcription, not a byte-identical recovery of a prior file.
% Owner, cohomology_audit and literature completed full proof/source reviews.
% No new numerical computation is claimed.

\begin{result}{E.218}{Proved}{A cyclic-Jacobian detector at every coefficient prime}
Let $k=\mathbb F_q$, let $\ell\ne\operatorname{char}k$ be a prime, and put
$U=\mathbb P^1_k\setminus\{0,1,\infty\}$. In the geometrically pro-$\ell$
arithmetic fundamental group of $U$, the section of an interior point of
$U(k)$ is not geometrically conjugate to any section carried by the
 decomposition group of a rational cusp.

There is an explicit characteristic finite geometric quotient detecting all
these distinctions. Put $g=\ell-1$ and
\[
 d=\begin{cases}[k(\mu_\ell):k],&\ell\text{ odd},\\
 [k(\mu_4):k],&\ell=2,\end{cases}
 \qquad Q=q^d.
 \tag{E.218.1}
\]
Choose an integer $N\ge1$ such that $\ell^N>2g(Q+1)$ and put
$c=(\ell-1)N-1$. If $\Pi=\pi_1(U_{\bar k})^{(\ell)}$, then the quotient
\[
 \Pi\big/\overline{\Pi^{\ell^N}\Gamma_{c+1}\Pi}
 \tag{E.218.2}
\]
with its full arithmetic extension over $G_k$ detects every interior-point/cusp
distinction above. In particular, an arithmetic construction that freezes an
interior reference Frobenius at a fixed off-diagonal place cannot retain it
through every finite pro-$\ell$ quotient while changing that place into
nonzero ordinary boundary support. This is a local obstruction to that
frozen-reference construction, not a finite-support theorem for arbitrary
global sections.
\end{result}

\begin{proof}
\emph{The family and its geometric monodromy.}
For $t\in U$, let $C_t$ be the smooth projective normalization of
\[
                  y^\ell=\frac{x(x-1)}{x-t}.
 \tag{E.218.3}
\]
The four branch points $0,1,t,\infty$ have orders $1,1,-1,-1$. They are
distinct on $U$, and the tame normalization defines a smooth proper family
over $U$. Riemann--Hurwitz gives
\[
 2g(C_t)-2=-2\ell+4(\ell-1),\qquad g(C_t)=\ell-1=g.
 \tag{E.218.4}
\]
Write $J_t=\operatorname{Jac}(C_t)$. Over a field containing $\mu_\ell$, the
deck transformation $\delta:y\mapsto\zeta_\ell y$ satisfies
$1+\delta+\cdots+\delta^{\ell-1}=0$ on $J_t$, since this norm is pullback
followed by pushforward through the quotient $C_t\to\mathbb P^1$, whose
Jacobian is zero. Thus
\[
 \mathcal O=\mathbb Z_\ell[\zeta_\ell],\qquad
 \lambda=1-\zeta_\ell
 \tag{E.218.5}
\]
acts on $T_\ell J_t$. The cyclotomic polynomial is irreducible over
$\mathbb Q_\ell$, so $T_\ell J_t$ is torsion-free over the discrete valuation
ring $\mathcal O$. Its $\mathbb Z_\ell$-rank is $2(\ell-1)$, hence it is free
of rank two over $\mathcal O$. In particular, $J_t[\lambda]$ has order
$\ell^2$.

Let $P_0,P_1,P_t,P_\infty$ denote the unique points over the four branch
points. The three branch differences $[P_0-P_\infty]$, $[P_1-P_\infty]$,
and $[P_t-P_\infty]$ are killed by both $\ell$ and $\lambda$. The divisor
of $y$ gives their relation
\[
 [P_0-P_\infty]+[P_1-P_\infty]-[P_t-P_\infty]=0.
 \tag{E.218.6}
\]
There are no further relations modulo $\ell$. Indeed, a principal divisor
supported on these branch points is invariant under $\delta$. Its defining
function $f$ therefore satisfies $\delta(f)/f\in\mu_\ell$, and consequently
$f=R(x)y^j$ for some $R(x)\in\bar k(x)^\times$. The branch valuations of
$R(x)$ are multiples of $\ell$, while those of $y^j$ are
$j(1,1,-1,-1)$. This proves the asserted uniqueness of the relation. The
branch differences therefore exhaust $J_t[\lambda]$.

They are sections over the parameter base, so geometric monodromy is
trivial on this group. Equivalently, on the free rank-two $\mathcal O$
Tate module its geometric image is contained in
\[
                1+\lambda\operatorname{Mat}_2(\mathcal O),
 \tag{E.218.7}
\]
a pro-$\ell$ group. The identification of the reduction modulo $\lambda$
with the branch-torsion module is used geometrically; the arithmetic
action may be semilinear on $\mathcal O$ and is retained throughout.

The actual arithmetic Tate representation consequently factors through
the geometrically pro-$\ell$ arithmetic fundamental group. More
explicitly, the kernel of the maximal pro-$\ell$ quotient of the geometric
fundamental group is characteristic, hence normal in the full arithmetic
group, and it acts trivially on this Tate module. This proves factorization
of the actual arithmetic representation, not merely a statement about its
geometric image.

\emph{The complete semistable model at the cusp.}
Work first at $t=0$. Blow up $(x,t)$ on $\mathbb P^1_{k[[t]]}$ and normalize
in the function field of E.218.3. Near the node of the blown-up base the
coordinates satisfy $xv=t$, and the normalized cover has equation
\[
 xv=t,\qquad y^\ell=\frac{x-1}{1-v}.
 \tag{E.218.8}
\]
At $x=v=0$ one has $y^\ell=-1$. Its derivative $\ell y^{\ell-1}$ is a unit,
so the normalized cover is etale over the base node there. All resulting
nodes therefore have thickness one, without any ramified extension of
the parameter field.

The two reduced components are
\[
 y^\ell=x-1,\qquad y^\ell=-\frac{z}{z-1},\quad x=tz.
 \tag{E.218.9}
\]
Both are projective lines: $y$ is a parameter on the first, and
$z=y^\ell/(y^\ell+1)$ on the second. They meet transversely at the $\ell$
points $y^\ell=-1$. Their generic vertical valuations are zero and their
residue equations are nonconstant separable covers, so both multiplicities
are one. At the remaining horizontal branch points the branch orders are
$1$ or $-1$; the usual local tame normalization is regular there. Away from
these branch points and the displayed nodes the cover is etale over a
regular chart. Thus these two components and their $\ell$ nodes are the
complete special fibre. Each component has self-intersection $-\ell$,
since the total reduced fibre has intersection zero with it. There is no
exceptional component of self-intersection $-1$ to contract.

For odd $\ell$ the node coordinates are $-\zeta_\ell^i$, and for $\ell=2$
they are the square roots of $-1$. The exact constant extension $k'/k$ in
E.218.1 therefore splits all components and nodes. The dual graph has
two vertices and $\ell$ edges, so its first Betti number is $\ell-1=g$.
The normalized components have genus zero. This model is semistable;
when $\ell=2$ it is the regular two-gon and is not asserted to be stable.
In that case the substitution $y'=y(x-t)$ identifies the family with the
Legendre equation $y'^2=x(x-1)(x-t)$.

The arithmetic semistable Tate filtration, including its full residue-field
action, now gives
\[
 0\longrightarrow\mathbb Z_\ell(1)^g
 \longrightarrow T_\ell J
 \longrightarrow\mathbb Z_\ell^g\longrightarrow0
 \tag{E.218.10}
\]
over $k'((t))$. The middle abelian-variety contribution of the general
filtration is zero because both normalized components have genus zero.
Here one may use \citet[author pp.~1--2, equations (1.1)--(1.2),
Theorem~1.5 and Corollary~1.6]{AuditDDM2021}; the stated local-field
and Henselian-DVR scope includes equal characteristic. Inertia acts
trivially on the two displayed graded pieces.

\emph{Arithmetic Frobenius and trace separation.}
Our Frobenius is the arithmetic automorphism $z\mapsto z^Q$ of $\bar k$
over $k'$. On the covariant Tate module it acts on $\mathbb Z_\ell(1)$ by
$Q$; for an abelian variety over $k'$ it is the action induced by the
$Q$-power Frobenius endomorphism. This convention, rather than the
contragredient convention on etale cohomology, is used in all the traces
below.

Every section carried by the cusp decomposition group sends arithmetic
Frobenius to a lift of this same residue-field Frobenius. Changing that
lift by any cusp-inertia element changes only the extension block in
E.218.10. Consequently its characteristic polynomial and trace on the
Tate module are
\[
 (X-Q)^g(X-1)^g,\qquad g(Q+1).
 \tag{E.218.11}
\]
This includes every tangential inertia parameter; no preferred tangent
has been imposed.

At an interior point $a\in U(k)$ the pullback representation is the actual
Tate module of the smooth abelian variety $J_a$. Its Frobenius trace
$A_a(Q)$ is an integer, and the Weil bound gives
\[
 |A_a(Q)|\le2g\sqrt Q.
 \tag{E.218.12}
\]
For the integrality, Frobenius/Tate identification, and the absolute
values of the eigenvalues, see \citet[Chapter~II, Theorem~1.1(b), printed
pp.~75--78]{AuditMilneAV2008}. Since $Q>1$,
\[
 0<g(Q+1)-A_a(Q)<2g(Q+1)<\ell^N.
 \tag{E.218.13}
\]
The traces are therefore different even modulo $\ell^N$. Geometrically
conjugate sections would induce conjugate representations and equal
traces, a contradiction. If the original sections over $k$ were conjugate,
their restrictions to $k'$ would also be conjugate. No injectivity assertion
for restriction in nonabelian cohomology is needed.

For either of the other two cusps, pull back this same family by an
automorphism of $U$ carrying that cusp to zero. Interior points remain
interior points. The same computation and the same numerical bound
apply. We do not assert that the original, untransformed family has the
same reduction type at all three cusps.

\emph{The finite nilpotent bound.}
Set $e=\ell-1$. The cyclotomic uniformizer satisfies
$\ell\mathcal O=\lambda^e\mathcal O$. Put
$K_a=1+\lambda^a\operatorname{Mat}_2(\mathcal O)$ for $a\ge1$.
Direct multiplication gives
\[
 [K_a,K_b]\subset K_{a+b}.
 \tag{E.218.14}
\]
The binomial formula gives $K_a^\ell\subset K_{a+e}$: for $1\le i<\ell$,
the coefficient $\binom\ell i$ has $\lambda$-valuation $e$, and the last
term has valuation $\ell a\ge a+e$. These statements include $\ell=2$.
Induction shows that every element of $K_1$ has its $\ell^N$-th power in
$K_{1+Ne}$, and that its $j$-th lower central term lies in $K_j$. Thus the
geometric image modulo $\ell^N$, which is modulo $\lambda^{eN}$, has
exponent dividing $\ell^N$ and nilpotence class at most $eN-1=c$.
It factors through E.218.2.

That quotient is finite: it is a finitely generated nilpotent pro-$\ell$
group of bounded exponent, so each of its finitely many lower-central
factors is a finitely generated abelian group of bounded exponent.
Its kernel is characteristic and hence stable under the full arithmetic
action. If two sections agreed up to geometric conjugacy in the resulting
arithmetic quotient, their Tate representations modulo $\ell^N$ would
agree up to conjugacy, contradicting E.218.13. The same characteristic
quotient works for all interior points and all three cusps.

Finally, an interior reference point at a fixed place and a point
specializing to a cusp yield the two distinct Frobenius types just computed.
In the ordinary full-section application, A.31 supplies the actual
off-diagonal local point class at nonzero support. Therefore retaining
the interior reference Frobenius through every finite quotient cannot
supply such an all-depth ordinary support match. No claim about the
number of ramified primes of a general global section, or about the
existence of a compatible global tower, follows.
\end{proof}
